\documentclass[../../../include/open-logic-section]{subfiles}

\begin{document}

\olfileid[id]{sth}{cardinals}{classing}
\olsection{Hingga, \usetoken{S}{enumerable}, \usetoken{S}{nonenumerable}}

Setelah mengenal kardinal, ada baiknya kita meluangkan sedikit
waktu untuk membahas berbagai jenis kardinal; khususnya, kardinal hingga,
!!{enumerable}, dan !!{nonenumerable}.

Dua hasil pertama kita menyiratkan bahwa kardinal hingga tepat merupakan
ordinal hingga, yang kita definisikan sebagai \emph{bilangan asli} dalam
\olref[z][infinity-again]{defnomega}:

\begin{prop}\ollabel{finitecardisoequal}
Misalkan $n, m \in \omega$. Maka $n = m$ jika dan hanya jika
$\cardeq{n}{m}$.
\end{prop}

\begin{proof}
Arah \emph{kiri-ke-kanan} sepele. Untuk membuktikan arah
\emph{kanan-ke-kiri}, andaikan $\cardeq{n}{m}$ meskipun $n \neq m$.
Berdasarkan Trikotomi, berlaku $n \in m$ atau $m \in n$; tanpa mengurangi
keumuman, andaikan $n \in m$. Maka $n \subsetneq m$ dan terdapat
!!a{bijection} $f \colon m \to n$, sehingga $m$ bersifat tak hingga
Dedekind, bertentangan dengan
\olref[z][infinity-again]{naturalnumbersarentinfinite}.
\end{proof}

\begin{cor}\ollabel{naturalsarecardinals}
Jika $n \in \omega$, maka $n$ merupakan suatu kardinal.
\end{cor}

\begin{proof}
Langsung.
\end{proof}
\noindent
Selain itu, beberapa gagasan yang wajar tentang arti menyebut suatu kardinal
``hingga'' atau ``tak hingga'' ternyata bertepatan:
\begin{thm}\ollabel{generalinfinitycharacter}Untuk sebarang himpunan $A$, hal-hal berikut ekuivalen:
\begin{enumerate}
	\item\ollabel{card:notinomega} $\card{A} \notin \omega$, yakni
	$\card{A}$ bukan bilangan asli;
	\item\ollabel{card:omegaplus} $\omega \leq \card{A}$;
	\item\ollabel{card:infinite} $A$ bersifat tak hingga Dedekind.
\end{enumerate}
\end{thm}

\begin{proof}
Dari \olref[ord-arithmetic][using-addition]{ordinfinitycharacter},
\olref[cardinals][cardsasords]{lem:CardinalsBehaveRight}, dan
\olref{naturalsarecardinals}.
\end{proof}

Hal ini membenarkan \emph{definisi} berikut bagi beberapa gagasan yang kita
gunakan secara agak informal dalam \olref[sfr][][]{part}:

\begin{defn}\ollabel{defnfinite}
Kita mengatakan bahwa $A$ bersifat \emph{hingga} jika dan hanya jika
$\card{A}$ merupakan bilangan asli, yakni $\card{A} \in \omega$. Jika tidak,
kita mengatakan bahwa $A$ bersifat \emph{tak hingga}.
\end{defn}
\noindent
Namun, perhatikan bahwa definisi ini disajikan dengan latar belakang
$\ZFC$. Bagaimanapun, kita memerlukan Pengurutan Baik untuk menjamin bahwa
setiap himpunan mempunyai suatu kardinalitas. Bahkan, tanpa Pengurutan Baik,
dapat terdapat suatu himpunan yang bukan hingga dan bukan pula tak hingga
Dedekind. Kita akan kembali membahas persoalan semacam ini dalam
\olref[choice][]{chap}. Untuk sekarang, kita terus mengandalkan Pengurutan
Baik.

Sekarang mari beralih dari kardinal hingga ke kardinal tak hingga. Berikut
dua hal mendasar:

\begin{cor}\ollabel{omegaisacardinal}
$\omega$ merupakan kardinal tak hingga terkecil.
\end{cor}

\begin{proof}
$\omega$ merupakan suatu kardinal, sebab $\omega$ bersifat tak hingga
Dedekind dan, jika $\cardeq{\omega}{n}$ bagi suatu $n \in \omega$, maka $n$
akan bersifat tak hingga Dedekind, bertentangan dengan
\olref[z][infinity-again]{naturalnumbersarentinfinite}. Sekarang $\omega$
merupakan kardinal tak hingga terkecil berdasarkan definisi.
\end{proof}

\begin{cor}
Setiap kardinal tak hingga merupakan ordinal limit.
\end{cor}

\begin{proof}
Misalkan $\alpha$ suatu ordinal penerus yang tak hingga, sehingga
$\alpha = \beta \ordplus 1$ bagi suatu $\beta$. Berdasarkan
\olref{finitecardisoequal}, $\beta$ juga tak hingga, sehingga
$\cardeq{\beta}{\beta \ordplus 1}$ berdasarkan
\olref[ord-arithmetic][using-addition]{ordinfinitycharacter}. Sekarang
$\card{\beta} = \card{\beta\ordplus 1} = \card{\alpha}$ berdasarkan
\olref[cardinals][cardsasords]{lem:CardinalsBehaveRight}, sehingga
$\alpha \neq \card{\alpha}$.
\end{proof}

Sejak \olref[sfr][siz][enm-alt]{defn:enumerable}, kita telah mencatat bahwa
himpunan tak hingga yang !!{enumerable} dapat dibedakan dari yang
!!{nonenumerable}. Definisi tersebut secara alami menghasilkan pernyataan
berikut:

\begin{prop}
$A$ bersifat !!{enumerable} jika dan hanya jika $\card{A} \leq \omega$, dan
$A$ bersifat !!{nonenumerable} jika dan hanya jika
$\omega < \card{A}$.
\end{prop}

\begin{proof}
Berdasarkan Trikotomi, kedua klaim ekuivalen, sehingga cukup dibuktikan bahwa
$A$ bersifat !!{enumerable} jika dan hanya jika
$\card{A} \leq \omega$. Untuk arah \emph{kanan-ke-kiri}: jika
$\card{A} \leq \omega$, maka $\cardle{A}{\omega}$ berdasarkan
\olref[cardinals][cardsasords]{lem:CardinalsBehaveRight} dan
\olref{omegaisacardinal}. Untuk arah \emph{kiri-ke-kanan}: andaikan $A$
bersifat !!{enumerable}; berdasarkan
\olref[sfr][siz][enm-alt]{defn:enumerable}, terdapat tiga kemungkinan:
\begin{enumerate}
	\item jika $A = \emptyset$, maka $\card{A} = 0 \in \omega$, berdasarkan
	\olref{naturalsarecardinals} dan
	\olref[cardinals][cardsasords]{lem:CardinalsBehaveRight}.
	\item jika $\cardeq{n}{A}$, maka $\card{A} = n \in \omega$, berdasarkan
	\olref{naturalsarecardinals} dan
	\olref[cardinals][cardsasords]{lem:CardinalsBehaveRight}.
	\item jika $\cardeq{\omega}{A}$, maka $\card{A} = \omega$, berdasarkan
	\olref{omegaisacardinal}.
\end{enumerate}
Jadi, dalam semua kasus, $\card{A} \leq \omega$.
\end{proof}
\noindent
Memang, $\omega$ menempati kedudukan khusus. Walaupun terdapat banyak ordinal
yang terhitung:

\begin{cor}
$\omega$ merupakan satu-satunya kardinal tak hingga yang !!{enumerable}.
\end{cor}

\begin{proof}
Misalkan $\cardfont{a}$ suatu kardinal tak hingga yang !!a{enumerable}.
Karena $\cardfont{a}$ tak hingga, $\omega \leq \cardfont{a}$. Karena
$\cardfont{a}$ merupakan kardinal yang !!a{enumerable},
$\cardfont{a} = \card{\cardfont{a}} \leq \omega$. Jadi
$\cardfont{a} = \omega$ berdasarkan Trikotomi. \end{proof}

Tentu saja, terdapat tak hingga banyak kardinal. Jadi, kita mungkin bertanya:
\emph{Ada berapa banyak kardinal?} Hasil-hasil berikut menunjukkan bahwa
kita mungkin perlu mempertimbangkan kembali pertanyaan itu.

\begin{prop}\ollabel{unioncardinalscardinal}
Jika setiap anggota $X$ merupakan suatu kardinal, maka $\bigcup X$ merupakan
suatu kardinal.
\end{prop}

\begin{proof}
Mudah diperiksa bahwa $\bigcup X$ merupakan suatu ordinal. Misalkan
$\alpha \in \bigcup X$ suatu ordinal; maka
$\alpha \in \cardfont{b} \in X$ bagi suatu kardinal $\cardfont{b}$. Karena
$\cardfont{b}$ suatu kardinal, $\cardless{\alpha}{\cardfont{b}}$. Karena
$\cardfont{b} \subseteq \bigcup X$, kita mempunyai
$\cardle{\cardfont{b}}{\bigcup X}$, sehingga
$\cardneq{\alpha}{\bigcup X}$. Dengan menggeneralisasi, $\bigcup X$
merupakan suatu kardinal.
\end{proof}

\begin{thm}\ollabel{lem:NoLargestCardinal}
Tidak terdapat kardinal terbesar.
\end{thm}

\begin{proof}
Untuk sebarang kardinal $\cardfont{a}$, Teorema Cantor
(\olref[sfr][siz][car]{thm:cantor}) dan
\olref[cardinals][cardsasords]{lem:CardinalsExist} menyiratkan bahwa
$\cardfont{a} < \card{\Pow{\cardfont{a}}}$.
\end{proof}

\begin{thm}
Himpunan semua kardinal tidak ada.
\end{thm}

\begin{proof}
Untuk memperoleh kontradiksi, andaikan
$C = \Setabs{\cardfont{a}}{\cardfont{a} \text{ merupakan suatu kardinal}}$.
Sekarang $\bigcup C$ merupakan suatu kardinal berdasarkan
\olref{unioncardinalscardinal}, sehingga berdasarkan
\olref{lem:NoLargestCardinal} terdapat suatu kardinal
$\cardfont{b} > \bigcup C$. Berdasarkan definisi $\cardfont{b} \in C$,
sehingga $\cardfont{b} \subseteq \bigcup{C}$, dan dengan demikian
$\cardfont{b} \leq \bigcup C$, suatu kontradiksi.
\end{proof}

Anda patut membandingkan hal ini dengan Paradoks Russell dan Burali--Forti.

\end{document}
